% ══════════════════════════════════════════════════════════════════════════
\section{Five-Dimensional Minimal Polytopes}
\label{sec:minimal-5d}

This section isolates the combinatorial input from
\citet{KreuzerSkarke1997} that controls which combined weight-system
patterns can occur in dimension~5.  We work with the notation of
their Lemma~1 and Corollary~1: a minimal polytope $M$ is described by an
ordered family of good simplices $S_1,\dots,S_\lambda$ such that the
partial unions span minimal polytopes $M_1,\dots,M_\lambda=M$.

\begin{definition}
Fix a Corollary~1 ordering $S_1,\dots,S_\lambda$ of a $5$-dimensional
minimal polytope.  Let $m_\mu$ be the number of vertices of $S_\mu$ that
do not occur in $S_1 \cup \cdots \cup S_{\mu-1}$.  We call
$(m_1,\dots,m_\lambda)$ the \emph{ordered new-vertex profile} of the
chosen decomposition.
\end{definition}

\begin{proposition}[Five-dimensional profile constraint]
\label{prop:5d-profiles}
Let $M$ be a $5$-dimensional minimal polytope and let
$(m_1,\dots,m_\lambda)$ be the ordered new-vertex profile of a Corollary~1
decomposition.  Then each $m_\mu \ge 2$ and
\[
  \sum_{\mu=1}^{\lambda} (m_\mu - 1) = 5.
\]
Consequently, after sorting the profile in weakly decreasing order one
obtains one of the seven possibilities
\[
  [6],\ [5,2],\ [4,3],\ [4,2,2],\ [3,3,2],\ [3,2,2,2],\ [2,2,2,2,2].
\]
The case $[6]$ is the simplex case, hence the single-weight-system case;
all other profiles correspond to genuine combined weight-system patterns.
\end{proposition}

\begin{proof}
Write $k_\mu$ for the number of vertices in $M_\mu$.  By
Corollary~1(a) of \citet{KreuzerSkarke1997},
$\dim M_\mu = k_\mu - \mu$.  Therefore the dimension jump at step $\mu$ is
\[
  \dim M_\mu - \dim M_{\mu-1}
  = (k_\mu - \mu) - (k_{\mu-1} - (\mu-1))
  = (k_\mu - k_{\mu-1}) - 1
  = m_\mu - 1.
\]
Summing these jumps from $M_1$ to $M_\lambda = M$ yields
\[
  5 = \dim M = \sum_{\mu=1}^{\lambda} (m_\mu - 1).
\]
Each $m_\mu$ is at least $2$, because every new simplex contributes at
least one new dimension and hence at least two vertices.  Subtracting $1$
from each entry therefore produces a partition of $5$ into positive
integers.  The only partitions of $5$ are
$5$, $4+1$, $3+2$, $3+1+1$, $2+2+1$, $2+1+1+1$, and $1+1+1+1+1$.
Adding $1$ back to each part gives exactly the seven profiles listed.
\end{proof}

\begin{lemma}[Lower-dimensional input data]
\label{lem:lower-dimensional-input}
Up to relabeling of vertices, the complete $3$-dimensional minimal-polytope
structures are
\[
  \{\{1,2,3,4\}\},\qquad
  \{\{1,2,3\},\{4,5\}\},\qquad
  \{\{1,2,3\},\{1,4,5\}\},\qquad
  \{\{1,2\},\{3,4\},\{5,6\}\},
\]
and the complete $4$-dimensional minimal-polytope structures are
\[
  \{\{1,2,3,4,5\}\},\qquad
  \{\{1,2,3,4\},\{1,2,5,6\}\},\qquad
  \{\{1,2,3,4\},\{1,5,6\}\},
\]
\[
  \{\{1,2,3,4\},\{5,6\}\},\qquad
  \{\{1,2,3\},\{4,5,6\}\},\qquad
  \{\{1,2,3\},\{4,5\},\{6,7\}\},
\]
\[
  \{\{1,2,3\},\{1,4,5\},\{6,7\}\},\qquad
  \{\{1,2,3\},\{1,4,5\},\{1,6,7\}\},\qquad
  \{\{1,2\},\{3,4\},\{5,6\},\{7,8\}\}.
\]
\end{lemma}

\begin{proof}
This is exactly the explicit $n=3$ and $n=4$ classification stated in
Section~4 of \citet{KreuzerSkarke1997}.  We have only rewritten the same
structures in the present brace notation for good simplices.
\end{proof}

\begin{lemma}[Dimension-$5$ extension rule]
\label{lem:5d-extension-rule}
Let $M$ be a $5$-dimensional minimal polytope.
\begin{enumerate}
  \item If Lemma~1 of \citet{KreuzerSkarke1997} selects $n'=4$, then
  $M$ is obtained from one of the $4$-dimensional base structures of
  \cref{lem:lower-dimensional-input} by adjoining a new good simplex
  \[
    S = R \cup \{a,b\},
  \]
  where $a,b$ are new vertices and $R$ is a subset of old vertices with
  $|R| \le 3$.
  \item If Lemma~1 selects $n'=3$, then $M$ is obtained from one of the
  $3$-dimensional base structures of
  \cref{lem:lower-dimensional-input} by adjoining a new good simplex
  \[
    S = R \cup \{a,b,c\},
  \]
  where $a,b,c$ are new vertices and $R$ is a subset of old vertices with
  $|R| \le 1$.
\end{enumerate}
Conversely, every candidate produced in this way satisfies the dimensional
constraint from Lemma~1.  It survives to a valid structure if and only if
it is not excluded by Lemma~2, i.e. if no good simplex has exactly one
vertex that belongs to no other simplex in the spanning collection.
\end{lemma}

\begin{proof}
Lemma~1 says that a non-simplex $5$-dimensional minimal polytope contains
an $n'$-dimensional minimal polytope $M'$ and one further good simplex $S$
such that
\[
  k-k' = 5-n'+1,
  \qquad
  \dim S \le n'.
\]
Since $n' \ge 5/2$, only $n'=4$ and $n'=3$ can occur.

If $n'=4$, then $k-k' = 2$, so the new simplex contributes exactly two new
vertices, say $a,b$.  Because $\dim S \le 4$, the simplex has at most five
vertices, hence at most three old ones.  Therefore
$S = R \cup \{a,b\}$ with $|R| \le 3$.

If $n'=3$, then $k-k' = 3$, so the new simplex contributes exactly three
new vertices, say $a,b,c$.  Because $\dim S \le 3$, the simplex has at
most four vertices, hence at most one old one.  Therefore
$S = R \cup \{a,b,c\}$ with $|R| \le 1$.

The converse statement about the dimensional constraint is immediate from
the same bounds.  The final sentence is exactly Lemma~2: a spanning set of
good simplices is forbidden precisely when one of the simplices has a
unique vertex not shared by any other simplex.
\end{proof}

\begin{theorem}[Complete $5$D structure list]
\label{thm:5d-structures}
Up to relabeling of the vertices, there are exactly $47$ overlap
structures for $5$-dimensional minimal polytopes.  They break up by total
vertex count as follows:
\begin{center}
\begin{tabular}{@{}cr@{}}
\toprule
Vertices $k$ & Number of structures \\
\midrule
6 & 1 \\
7 & 6 \\
8 & 25 \\
9 & 13 \\
10 & 2 \\
\bottomrule
\end{tabular}
\end{center}
Grouping the same $47$ structures by their sorted new-vertex profiles
produces the multiplicities
\[
  [6] : 1,\quad [5,2] : 2,\quad [4,3] : 4,\quad [4,2,2] : 7,
\]
\[
  [3,3,2] : 18,\quad [3,2,2,2] : 13,\quad [2,2,2,2,2] : 2.
\]
Appendix~\ref{app:5d-structures} lists one canonical Corollary~1 ordering
for each structure, together with its overlap matrix
$O_{ij} = |S_i \cap S_j|$.
\end{theorem}

\begin{proof}
By \cref{lem:lower-dimensional-input}, there are exactly four possible
$3$-dimensional base structures and nine possible $4$-dimensional base
structures.  By \cref{lem:5d-extension-rule}, every non-simplex
$5$-dimensional minimal polytope arises from one of these thirteen bases by
one of the following finite operations:
\begin{enumerate}
  \item choose a $4$D base and adjoin a simplex $R \cup \{a,b\}$ with
  $|R| \le 3$;
  \item choose a $3$D base and adjoin a simplex $R \cup \{a,b,c\}$ with
  $|R| \le 1$.
\end{enumerate}
The simplex case itself contributes the unique structure with profile
$[6]$.

Thus the problem has been reduced to a completely finite search.  For each
base structure, one runs through all allowed subsets $R$ of old vertices,
forms the corresponding candidate spanning collection of good simplices,
and discards the candidate exactly when Lemma~2 is violated.  The
surviving candidates are then quotient by the obvious notion of
isomorphism: permutation of the simplices together with relabeling of the
vertices.  Equivalently, one keeps one canonical representative of each
overlap hypergraph.

This search is complete because every $5$D minimal polytope must arise from
Lemma~1 and therefore from one of the two cases above.  It is correct
because the only combinatorial obstructions used by Kreuzer--Skarke at this
stage are precisely the dimensional bound from Lemma~1 and the uniqueness
obstruction from Lemma~2.  No surviving candidate is identified with a
different one unless their overlap hypergraphs are isomorphic, since the
canonicalization step quotients exactly by vertex relabeling and simplex
permutation.

Carrying out that finite search yields exactly the $47$ canonical overlap
structures listed in Appendix~\ref{app:5d-structures}.  Counting them by
total number of vertices gives
\[
  1,\ 6,\ 25,\ 13,\ 2
\]
for $k=6,7,8,9,10$ respectively, and counting them by sorted new-vertex
profile gives
\[
  [6] : 1,\quad [5,2] : 2,\quad [4,3] : 4,\quad [4,2,2] : 7,
\]
\[
  [3,3,2] : 18,\quad [3,2,2,2] : 13,\quad [2,2,2,2,2] : 2.
\]
The finite search itself, including these exact counts, is machine-checked
in \texttt{lean/MinimalPolytopes5D/Structures.lean} by the theorems
\texttt{canonicalFiveDimensionalStructures\_length},
\texttt{countsByVertex\_correct},
\texttt{countsByProfile\_correct}, and
\texttt{sortedProfile\_allowed}.
\end{proof}

\begin{remark}
The appendix records one admissible Corollary~1 ordering for each overlap
hypergraph.  Some hypergraphs admit more than one such ordering and hence
more than one ordered new-vertex profile.  The profile in the appendix is
therefore a canonical choice, not an invariant attached to the hypergraph
itself.  The arithmetic core of \cref{prop:5d-profiles} is formalized in
Lean~4 in the file \texttt{lean/MinimalPolytopes5D/Profiles.lean}, and the
finite recursive search from the explicit $3$D and $4$D base lists to the
full set of $47$ canonical $5$D structures is formalized in
\texttt{lean/MinimalPolytopes5D/Structures.lean}.
\end{remark}